\begin{definition}
Let $\Omega\subset\mathbb{C}$, let $j\in\{0,1\}$ and let
$u,v\colon\mathbb{C}\to\mathbb{C}$.  Identify $\mathbb{C}$ with $\mathbb{R}^2$
by $z=x+iy$ and write $\partial_0=\partial/\partial x$,
$\partial_1=\partial/\partial y$.  We say that $v$ is the
\emph{distributional $\partial_j$-derivative of $u$ on $\Omega$} if
\[
  \int_{\mathbb{C}}u(z)\,\partial_j\phi(z)\,d\sigma(z)
  =-\int_{\mathbb{C}}v(z)\,\phi(z)\,d\sigma(z)
\]
for every test function $\phi$ for $\Omega$ in the sense of
\noderef{IsTestFn}.  (Both integrands vanish off the compact support of
$\phi$, so both integrals are taken over a compact subset of $\Omega$.)

This is the usual meaning of the phrase ``distributional derivative
$D^\nu u$'' in Section~1 of the paper, for a multi-index $\nu$ of order one.
\end{definition}
